\section{Related Work}\label{sec:related}

\textbf{Certifiable differentiable layers.} SDPRLayers \cite{holmes2025sdpr} is the closest work: it lifts a homogeneous QCQP to a Shor SDP, reads off a global solution when the relaxation is tight, and differentiates the \emph{original} KKT system with an implicit-selection argument that tolerates redundant equality constraints. Three differences matter. (i)~\emph{Constraint class}: their theory is for quadratic \emph{equalities}; inequalities require slacks that inflate the PSD block and break strict complementarity, whereas ours are native. (ii)~\emph{Where the SDP lives}: they solve an SDP in the forward pass for every sample and step; we solve a fixed-dimension polynomial system forward and use an SDP only for the posteriori certificate, asynchronously. (iii)~\emph{What global optimality rests on}: theirs rests on the forward relaxation, ours on a separate $14\times14$ relaxation, so a failure of tightness costs the certificate but not the forward solve.

\textbf{Algebraically global quaternion solvers.} QPEP \cite{wu2022qpep} unifies PnP, PnL, hand--eye calibration, point-to-plane registration, conic pose, range-based relative pose and Stewart-platform kinematics as ``quadratic pose estimation problems'', eliminates $\mathbf t$, and enumerates all KKT roots with a Gr\"obner basis (up to 40 candidates). Our work adds inequalities and gradients to QPEP: we keep its elimination and enumeration and add affine constraints, an ALM loop, FB differentiation and a semidefinite posteriori certificate. In QPEP's sense, global optimality means the best stationary point given a complete enumeration, which is not a convex certificate.

\textbf{Differentiable convex-optimisation layers.} OptNet \cite{amos2017optnet} differentiates QPs via the KKT system of an interior-point solve, CVXPYLayers \cite{agrawal2019cvxpylayers} does so for cone programs, and dQP \cite{magoon2025dqp} decouples solver and derivative and shows that the active set determines both. These methods assume convexity, an identified active set and LICQ or strict complementarity; our feasible set lives on $\mathbb{S}^3$, LICQ fails at four contacts and three out-of-plane degrees of freedom, and strict complementarity fails at grazing contact. We borrow dQP's idea of reducing the derivative to a small system fixed by the active set, and add a rank-robust regularisation plus a limit theory.

\textbf{Local pose layers and non-smooth differentiation.} BPnP \cite{chen2020bpn} back-propagates through a PnP solve by implicit differentiation, EPro-PnP \cite{chen2022epropnp} integrates a pose posterior by Monte Carlo, and Theseus \cite{pineda2022theseus} provides differentiable nonlinear least squares on manifolds with implicit or unrolled derivatives. All are local and support penalties but not hard constraints. The Fischer--Burmeister function \cite{fischer1992} and smoothed variants \cite{kanzow1996} turn complementarity into a semi-smooth equation \cite{facchinei2003}; the non-convex ALM is classical \cite{bertsekas1982,rockafellar1998}, with exact-penalty variants in \cite{fletcher1987}; and \cite{bolte2021nonsmooth} is the framework in which we state the limit of Prop.~\ref{prop:limit}. Table~\ref{tab:positioning} summarises the comparison in the appendix.
